\subsection{A critical sampler from a column envelope}
\label{sec:column-envelope}

The first-layer rank can grow with the width in the result below.  A
different condition replaces the rank bound.  Define the column envelope
of the first weight matrix by
\[
 C_{\mathrm{col}}=\sum_{j=1}^n\max_i|w_{1,ij}|.
\]
This quantity is computed directly from the weights.  It is one for a
Hadamard sign matrix divided by its width, although that matrix has full
rank.  It can also be as large as $n$.  The theorem therefore leaves the
unrestricted critical problem open.

The empirical telescope uses the multilevel debiasing principles of
Rhee and Glynn~\cite{rheeglynn2015} and the antithetic vector-mean
construction of Blanchet and Glynn~\cite{blanchetglynn2015}.
We give explicit bounded correction branches, sampled derivatives,
and a finite-bit implementation for the present network.

\begin{theorem}[Column-envelope sampling]
\label{thm:column-envelope-critical}
Put $K_{\mathrm{col}}=\max\{1,C_{\mathrm{col}}\}$.
Consider zero-bias tanh networks with row sums at most one.  At
$D=\lceil\log_2 n\rceil$ and $b=O(\log n)$, weight-only preprocessing
and storage $O(Dn^2\log^6 n)$ suffice for an exact sampler with
\[
 \mathbb E Q=O(K_{\mathrm{col}}^2\sqrt{D+1}),\qquad
 \mathbb E W=O(K_{\mathrm{col}}^7\sqrt{D+1}\,B^{40}),
 \quad B=16+b+\lceil\log_2(nD+2)\rceil.
\]
Here $Q$ counts input probes and $W$ counts all online bit operations.
The constants in the online bounds are absolute.  The preprocessing
bound uses the stated precision convention $b=O(\log n)$; the
additional preprocessing in this proof is explicitly $O(Dn^2B^4)$.
The worst-case probe order on $C_{\mathrm{col}}\le C_0$ for every
fixed $C_0\ge1$ is
$\Theta(\sqrt D)$ under the width and grid conditions of the critical
mean-chain lower bound.
\end{theorem}

\paragraph{Proof idea.}
One context probe supplies an implicit bounded vector whose expectation
is the entire first preactivation, up to a known scale.  A corrected
empirical-mean factory needs $O(K_{\mathrm{col}}^2D)$ such vectors
on its open amplitude branch.  Its correction terms involve fourth derivatives.
We sample their chain-rule expansions along at most four paths through
the network.  A small amount of slack in each first-derivative factor
keeps their product bounded by an absolute constant.  Fourth-order
antithetic remainders then admit an exact importance sampler.  Every
subsequent neuron coin uses the already probed empirical observations;
no dense evaluation is performed.

The proof occupies the rest of this subsection. Its first part builds
the source vector, the amplitude gate, and the common derivative
majorants. The second gives the coin for one chain-rule diagram, the
third the empirical correction and its antithetic identity, and the
fourth the exact tilting that samples the fourth-order remainders. The
last part assembles the telescope and counts sources and bit work.

\subsubsection{A source vector and common derivative bounds}

For a nonzero first-layer column put $d_j=\max_i|w_{1,ij}|$, and let
$c_j$ be the least power of two at least $d_j$.  Set $c_j=0$ on a zero
column and $C=\sum_jc_j$.  Then
\[
 C_{\mathrm{col}}\le C\le2C_{\mathrm{col}}.
\]
The zero matrix has fair output and is handled separately.  Otherwise
choose $J$ with probability $c_J/C$, query $x_J$, and retain the pair
$(J,x_J)$.  It represents the vector
\begin{equation}
 Y_i=w_{1,iJ}x_J/c_J,\qquad |Y_i|\le1,
 \qquad \mathbb E Y=W_1x/C.
 \label{eq:column-vector-source}
\end{equation}
Only requested coordinates are formed.  Each is a known dyadic number.
Indeed $2^{-b}\le c_J\le1$ is a power of two, so its denominator divides
the original grid denominator $2^b$.  Averaging $m$ such observations
therefore introduces only $O(\log m)$ additional bits.

Let $F(y)$ be the scalar output obtained by first applying
$\tanh(Cy)$ coordinatewise and then applying layers $2,\ldots,D$ of
the given network.  Thus $F(\mathbb E Y)$ is the desired output mean.
Put $r_1=1$.  For $\ell\ge2$, choose dyadic radii with
\[
 \tanh r_{\ell-1}\le r_\ell
          \le\tanh r_{\ell-1}+\epsilon_0,
 \qquad 0<\epsilon_0\le(16D^2)^{-1}.
\]
Certified scalar evaluation and upward rounding construct these with
$O(\log(D+2))$ bits.  The one-Lipschitz property of tanh gives
\[
 r_\ell\le R_{\ell-1}+D\epsilon_0,
 \qquad R_j=\tanh^{\circ j}(1),
 \qquad r_\ell^2\le2/\ell,\qquad r_\ell\le1.
\]
For the square bound, expand
$(\sqrt{3/(2\ell+1)}+1/(16D))^2$ and use $\ell\le D$; its
three terms sum to less than $2/\ell$.  The upper bound one follows
inductively from $\tanh1+1/16<1$.  Also
$r_D\ge R_{D-1}\ge(2D-1)^{-1/2}$, using the scalar radius lower
bound.  We work with $D\ge32$; the finitely many smaller depths use
the existing universal critical sampler, whose cost is absorbed by the
absolute constants.  Let $q$ be the least power of two at
least $2r_D$.  Then
\begin{equation}
 0<q\le1,\qquad q\le6D^{-1/2},\qquad q^{-1}\le\sqrt D,
 \qquad |f|\le\tfrac12,\quad f=F/q.
 \label{eq:column-top-gate}
\end{equation}

A \emph{common derivative majorant} is a nonnegative symmetric tensor
$C^{(k)}$ such that
$|\partial_{i_1}\cdots\partial_{i_k}f(y)|\le C^{(k)}_{i_1\ldots i_k}$
for every $y\in[-1,1]^n$.  Its mass is $M_k=\sum C^{(k)}$.
The tensor is independent of the point $y$.  This is stronger than a
pointwise bound on the sum of derivative magnitudes.

The positive chain-rule recursion, with scalar bounds
\begin{equation}
 c_{\ell,1}=1,\quad c_{\ell,2}=2r_\ell,
 \quad c_{\ell,3}=2,\quad c_{\ell,4}=16r_\ell,
 \label{eq:column-scalar-majorants}
\end{equation}
constructs common majorants through order four.  To check their masses,
divide order $k$ by $C^k$ and let $A_\ell,B_\ell,E_\ell,H_\ell$
denote the resulting masses at orders one through four.  The usual
chain rule gives
\begin{align*}
 A_\ell&\le1,&
 B_\ell&\le B_{\ell-1}+2r_\ell,\\
 E_\ell&\le E_{\ell-1}+6r_\ell B_{\ell-1}+2,&
 H_\ell&\le H_{\ell-1}+8r_\ell E_{\ell-1}
       +6r_\ell B_{\ell-1}^2+12B_{\ell-1}+16r_\ell.
\end{align*}
Consequently
\[
 B_D\le4\sqrt2\sqrt D<6\sqrt D,\quad E_D\le50D,
 \quad H_D\le928D^{3/2}.
\]
For the last inequality, the increment is at most
$656\sqrt2\sqrt\ell<928\sqrt\ell$.
Writing $\kappa=\max\{1,C\}$, sufficient rational upper bounds for
$f=F/q$ are
\begin{equation}
 M_2\le6\kappa^2D,
 \qquad M_3\le50\kappa^4D^2,
 \qquad M_4\le928\kappa^4D^2.
 \label{eq:column-derivative-masses}
\end{equation}

The full tensors need not be stored.  At each neuron and each order
$k\le4$, store the scalar mass obtained by putting one in every input
direction.  A chain-rule term is indexed by a partition of $k$ labeled
directions.  Its weight is $c_{\ell,|\pi|}$ times the product of the
corresponding lower-order row masses.  Prefix arrays select both a
partition and each predecessor.  This samples an ordered chain-rule
diagram with the correct normalized majorant weight.  At most $k$
paths cross any layer, and at most $k-1$ vertices have derivative
order greater than one.  Its leaves supply the ordered tensor indices. The root factor $1/q$
scales the tensor and its mass together, so it does not change the
diagram distribution. Throughout, $M_k$ denotes the actual mass of
the constructed tensor; the displayed inequalities are upper bounds
used only in the work estimates.
All stored quantities are dyadic, with $O(D[b+\log(D+n+2)])$ bits.
The four mass passes and their selectors use $O(Dn^2)$ scalar
operations and storage entries.

\subsubsection{A bounded coin for a chain-rule diagram}

The normalized-neuron sampler used below has the following interface.
Its first layer uses the arbitrary-radius majority factory of
Theorem~\ref{thm:quadraticradiusbits}, with radius $C$ and its outer
gate $\tanh C$.  The resulting mean is $\tanh(Cy_i)$, which already
equals $h_{1,i}/r_1$ because $r_1=1$.
At a later row put $\lambda=\sum_j|w_{\ell,ij}|$ and
$\rho=\lambda r_{\ell-1}$.  If $\lambda=0$, return a fair sign.
Otherwise select signed normalized predecessor coins with probabilities
$|w_{\ell,ij}|/\lambda$.  Their mean $z$ satisfies
$\rho z=\sum_jw_{\ell,ij}h_{\ell-1,j}$.
The scalar majority factory has mean
$\tanh(\rho z)/\tanh\rho$; a final known gate
$\tanh\rho/r_\ell\le1$ gives mean $h_{\ell,i}/r_\ell$.
The inequality follows from $\rho\le r_{\ell-1}$ and the upward
radius bound.  A further gate of probability $r_\ell$ gives an
actual-activation sign of mean $h_{\ell,i}$ when required.

\begin{lemma}[A sampled derivative term]
\label{lem:column-diagram-coin}
Fix $k\le4$, a sampled order-$k$ diagram, and an input point supplied
by coordinate-coin access.  Let $\alpha$ be the product of its edge
signs and its scalar derivative factors divided by the
majorants in~\eqref{eq:column-scalar-majorants}.  There is an exact
$Z\in\{-1,0,1\}$ with
\[
 \mathbb E Z=\alpha/48.
\]
It makes a polynomial number of coordinate-source calls in expectation.
The polynomial is uniform over the supplied point and over all diagrams.
\end{lemma}

\paragraph{Proof idea.}
The sharp bound $\tanh'\le1$ has an endpoint that prevents a direct
positive factory.  Divide each such factor by $1+1/(4D)$ instead.
Their total loss stays below three.  Every higher derivative is a
fixed polynomial in a normalized activation, and there are at most
three such vertices in a fourth-order diagram.

\begin{proof}
At a vertex in layer $\ell$, put $h=\tanh s$, $u=h/r_\ell$.
Independent normalized-neuron coins have mean $u$.  The polynomials
for derivative orders two, three, and four, after division by their
majorants, are respectively
\[
 -u+r_\ell^2u^3,\qquad
 -1+4r_\ell^2u^2-3r_\ell^4u^4,\qquad
 u-\tfrac52r_\ell^2u^3+\tfrac32r_\ell^4u^5.
\]
Their absolute coefficient sums are at most $2,8,5$.  Sampling a
signed monomial from these coefficients, and returning a fair sign
on unused mass, gives a sign with the corresponding mean divided by
$K_2=2,K_3=8,K_4=5$.

For a first-derivative vertex put $\epsilon=1/(4D)$ and $M=4D+1$.
Take $M$ independent signs of mean $h$, and let $K$ be the number of
plus signs.  Return a bit with conditional probability
\begin{equation}
 \frac{4K(M-K)}{M(M-1)(1+\epsilon)}.
 \label{eq:column-first-derivative-factory}
\end{equation}
This belongs to $[0,1]$, because
$4K(M-K)\le M^2$ and $M/(M-1)=1+\epsilon$.
Its expectation is $(1-h^2)/(1+\epsilon)$.

Use fresh neuron coins in all these constructions.  There are at most
$4D$ first-derivative vertices, so their lost factor is below
$(1+1/(4D))^{4D}<3$.  The identity
$\sum(\nu-1)=k-1\le3$ bounds the product of the higher-derivative
losses by $16$: the largest possibility is one order-three and one
order-two vertex.  Multiply all sampled signs and bits and the known
edge signs.  Finally retain that product with probability
\[
 (1+\epsilon)^{N_1}\prod_{\nu\ge2}K_\nu/48,
\]
returning zero otherwise.  This is a known rational probability at
most one.  Conditional independence proves the required mean.
The work bound is proved quantitatively below.
\end{proof}

Conditional on a majorant diagram with leaf indices $J_1,\ldots,J_k$,
the resulting bounded estimate of a directional derivative is
\begin{equation}
 48M_k Z\prod_{a=1}^k v_{a,J_a}.
 \label{eq:column-directional-estimator}
\end{equation}
Its expectation is $D^kf(y)[v_1,\ldots,v_k]$.  Diagram sampling
requires only the scalar masses and prefix arrays already stored.

\subsubsection{An empirical correction with sparse coordinate access}

For $m\ge2$ vector observations let $z$ be their mean, $V_m$ their
unbiased sample covariance, and define
\[
 G_m=f(z)-\frac{V_m:D^2f(z)}{2m},\qquad P_m=\mathbb E G_m.
\]
The colon denotes contraction of the two indices.  The strong law
and bounded convergence give $\mathbb E f(z)\to f(\mathbb E Y)$.
Every covariance entry has magnitude at most two, so the correction
has magnitude at most $M_2/m$ and $P_m\to f(\mathbb E Y)$.
If $a,b$ are a uniform unordered
pair of observations and $u=Y_a-Y_b$, then
$\mathbb E[uu^\top]/2=V_m$.  Thus the base coefficient has the bounded
unbiased estimator
\[
 \widehat G_m=\tfrac12 S-\frac{48M_2}{4m}
                         Z\,u_{J_1}u_{J_2},
 \qquad \mathbb E S=2f(z).
\]
The sign $S$ is supplied by the normalized critical sampler followed
by its gate $2r_D/q\le1$.  At an empirical point $z$, each coordinate
source call simply selects one retained observation uniformly and
samples its known dyadic coordinate.  The second term uses
Lemma~\ref{lem:column-diagram-coin}.  If $m\ge384M_2$, then
$|\widehat G_m|\le5/8$.

We next give an exact correction for a fixed population of $2m$
observations.  Let its mean be $z$.  A uniform half $A$ and its
complement $B$ have means $z+\delta,z-\delta$ and covariance matrices
$V_A,V_B$.  Write $H=D^2f(z)$ and $T=D^3f(z)$.  Define
\begin{align*}
 R_4&=\tfrac12[f(z+\delta)+f(z-\delta)]-f(z)
                          -\tfrac12H[\delta,\delta],\\
 R_H&=\tfrac12V_A:[D^2f(z+\delta)-H-T[\delta]]
       +\tfrac12V_B:[D^2f(z-\delta)-H+T[\delta]].
\end{align*}
The exact identity is
\begin{equation}
 G_{2m}-\mathbb E_A G_m(A)
 =c_3+\mathbb E_A\left[\frac{R_H}{2m}-R_4\right],
 \quad
 c_3=\frac{D^3f(z):\sum_{a=1}^{2m}(Y_a-z)^{\otimes3}}
                  {4m(m-1)(2m-1)}.
 \label{eq:column-antithetic-identity}
\end{equation}
Here $T[\delta]$ leaves two indices uncontracted.

\paragraph{Proof of the identity.}
The within-population covariance decomposition and complement symmetry
give
\[
 (m-1)(V_A+V_B)=(2m-1)V_{2m}-2m\delta\delta^\top,
 \quad \mathbb E\delta\delta^\top=V_{2m}/(2m),
 \quad \mathbb E\delta^{\otimes3}=0.
\]
For centered observations $v_a=Y_a-z$, subset inclusion probabilities
$1/2$ and $(m-1)/(2(2m-1))$ also give
\[
 \mathbb E[V_A\otimes\delta]
 =\frac{\sum_av_a^{\otimes3}}{2(m-1)(2m-1)}.
\]
Indeed the distinct-index terms sum to minus the equal-index third
moment, while the third moment of $\delta$ vanishes.  Substituting
these expressions in the paired Taylor expansions proves
\eqref{eq:column-antithetic-identity}.

The cubic term is estimated by a uniform observation $a$, the
direction $u=Y_a-z$, and an order-three diagram.  Its estimator is
\[
 \frac{48M_3 Z\,u_{J_1}u_{J_2}u_{J_3}}
                         {2(m-1)(2m-1)}.
\]
Its magnitude is at most $48\cdot6M_3/m^2$.

For the fourth-order terms use integral remainders.  With a fair sign
$\sigma$, an unordered pair $u=Y_a-Y_b$ from the half of mean
$z+\sigma\delta$, and independent beta variables, they are
\begin{align}
 R_4&=\frac1{24}\mathbb E_{\sigma,t\sim\operatorname{Beta}(1,4)}
    D^4f(z+\sigma t\delta)[\delta,\delta,\delta,\delta],
 \label{eq:column-quartic-integral}\\
 \mathbb E_A\frac{R_H}{2m}
 &=\frac1{8m}\mathbb E_{A,\sigma,u,t\sim\operatorname{Beta}(1,2)}
    D^4f(z+\sigma t\delta)[u,u,\delta,\delta].
 \label{eq:column-hessian-integral}
\end{align}
The first formula is the fourth-order Taylor remainder with density
$4(1-t)^3$; the second is the second-order remainder of the Hessian
with density $2(1-t)$.  These formulas avoid subtracting nearly equal
random function estimates.

\subsubsection{Exact tilting of a half-sample}

Choose a uniform pairing of the $2m$ observations and select one
member of each pair by independent fair signs $\varepsilon_j$.
For any fixed coordinate $i$,
\[
 \delta_i=\sum_{j=1}^m\varepsilon_jd_j,
 \quad d_j=(Y_{a_j,i}-Y_{b_j,i})/(2m),\quad |d_j|\le1/m.
\]
For a fixed pairing put $v=\sum d_j^2$, $u=\sum d_j^4$.  Its exact
second and fourth moments are
\begin{equation}
 \nu_2=v\le1/m,\qquad
 \nu_4=3v^2-2u\le3/m^2.
 \label{eq:column-pairing-moments}
\end{equation}
To sample signs with density proportional to
$(\sum\varepsilon_jd_j)^k$, $k=2,4$, expose them sequentially.
At partial sum $s$ the remaining moment polynomials are
\[
 P_2(s)=s^2+v,\qquad
 P_4(s)=s^4+6s^2v+3v^2-2u.
\]
The next plus probability is $P_k(s+d)/(2P_k(s))$, using the
remaining moments after this increment in the numerator.  The two
probabilities add to one.  A zero denominator is a null path under
the tilted law.  This is a finite rational Doob transform.

Accept a fresh pairing with probability $m\nu_2$ for a quadratic
tilt, or $m^2\nu_4/3$ for a quartic tilt.  On acceptance use its
tilted signs.  On rejection return zero.  There is no retry loop.
Consequently the accepted half-sample has unnormalized density
$m\delta_i^2$ or $m^2\delta_i^4/3$ relative to a uniform half.

Here is the complete Hessian-remainder branch.  Its mass is
$48M_4/(2m^2)$.  Sample an order-four diagram and then choose
$i=J_3$ or $J_4$ equally.  Perform the quadratic tilt in coordinate
$i$.  Put $v_\delta=(\delta_{J_3}^2+\delta_{J_4}^2)/2$.
On an accepted split, generate $\sigma,u,t$ as in
\eqref{eq:column-hessian-integral}, and the diagram coin $Z$ at
$z+\sigma t\delta$.  Return a bounded sign estimator with conditional
mean
\[
 Z\,\frac{u_{J_1}u_{J_2}}4
      \frac{\delta_{J_3}\delta_{J_4}}{v_\delta}.
\]
Both rational factors belong to $[-1,1]$; accepted splits have
$v_\delta>0$.  Multiplying the branch mass, the accepted density
$mv_\delta$, and the diagram mean $\alpha/48$ gives exactly the
integrand coefficient $M_4/(8m)$ in
\eqref{eq:column-hessian-integral}.

The quartic branch has mass $48M_4/(8m^2)$.  Sample an order-four
diagram, choose one of its four leaves uniformly, and perform the
quartic tilt in that coordinate.  Put
$w_\delta=\frac14\sum_{a=1}^4\delta_{J_a}^4$.
On acceptance generate $\sigma,t$ from
\eqref{eq:column-quartic-integral} and return conditional mean
\[
 -Z\,\frac{\prod_{a=1}^4\delta_{J_a}}{w_\delta}.
\]
The ratio has magnitude at most one by the arithmetic-geometric mean
inequality.  The branch mass and accepted density
$m^2w_\delta/3$ give precisely the negative term in
\eqref{eq:column-quartic-integral}.  Only four coordinate sequences
are ever needed to form these moments and ratios.

The beta variables require no real-number oracle.  Use one P\'olya
urn, initially $(1,4)$ or $(1,2)$, for the whole integral term.  A
coordinate coin of mean
$z_i+\sigma t\delta_i=(1-t)z_i+t(z_i+\sigma\delta_i)$ first draws
from that urn, then selects a uniform observation from the full
population or the chosen half.  It finally samples the known dyadic
coordinate $Y_i$.  Conditional on the beta parameter, all urn draws
are independent.  The same urn is shared by every neuron and every
factor in that one diagram.  A fresh urn per factor would change the
integral and is not used. To verify this representation at finite
stopping times, a specified urn word with $s$ successes and $f$
failures has probability
\[
 \frac{(a)_s(b)_f}{(a+b)_{s+f}}
 =\frac{1}{\mathrm B(a,b)}\int_0^1
       t^{s+a-1}(1-t)^{f+b-1}\,dt.
\]
Here $(a)_s$ is a rising factorial. The first equality follows by
multiplying the predictive probabilities; the second follows by
integrating by parts. Hence every finite semantic transcript of source
outputs and high-level choices has exactly the beta-mixture probability,
even when requests are adaptive. The
finite moment bounds below imply almost sure termination, so the
same identity applies to the returned value.

\subsubsection{Exactness, source count, and finite-bit work}

Put $A=6M_3+M_4$ and $K_*=48$.  Let $m_0$ be the least power of two with
\begin{equation}
 m_0\ge2,\qquad m_0\ge8K_*M_2,
 \qquad m_0^2\ge16K_*A.
 \label{eq:column-starting-degree}
\end{equation}
Every bounded rational estimator below is rounded to a sign of that
conditional mean using a fresh known coin.  Give the base mass $3/4$,
and return a known sign of mean
$\widehat G_{m_0}/(3/4)$.  At each level $m=2^am_0$, give the
cubic, Hessian, and quartic branches their masses just specified.
Their sum is at most $K_*A/m^2$.  Return a fair sign on all remaining
mass, including rejected tilts.  The total non-base mass is at most
$4K_*A/(3m_0^2)\le1/12$.  This countable choice is exact with finite
rational arithmetic: its actual non-base mass is
$4(288M_3+30M_4)/(3m_0^2)$.  Conditional on that choice, take level
$a$ with probability $(3/4)4^{-a}$, using a geometric draw, and
choose the three branches in the proportions
$288M_3:24M_4:6M_4$.  Zero masses are omitted.
Every selected branch is bounded and its
mean is its required coefficient.  Therefore
\[
 P_{m_0}+\sum_{a\ge0}(P_{2^{a+1}m_0}-P_{2^am_0})
       =f(\mathbb E Y)
\]
proves exactness.  Absolute convergence follows from the branch-mass
sum.  If a majorant mass is zero, its branch is omitted.

The base probes $m_0$ observations and a correction probes $2m$.
All later neuron-source requests select from those retained
observations.  Hence
\begin{equation}
 \mathbb E N\le\tfrac34m_0+
       \sum_{m=2^am_0}\frac{K_*A}{m^2}(2m)
       =\tfrac34m_0+\frac{4K_*A}{m_0}\le m_0.
 \label{eq:column-source-count}
\end{equation}
For explicit constants, the bounds
\eqref{eq:column-derivative-masses} give
$8\cdot48M_2\le2304\kappa^2D$ and
$16\cdot48(6M_3+M_4)\le943104\kappa^4D^2$.
A power of two at least $2304\kappa^2D$ therefore satisfies the
three starting conditions, so the least valid power obeys
$m_0<4608\kappa^2D$.  Opening the gate $q$ and otherwise returning
a fair sign proves
$\mathbb E Q<27648\kappa^2\sqrt D$ for $D\ge32$.

We finish the bit count, including the unbounded urn counters.
Use full majorities in every normalized neuron call.  Every monomial
and every first-derivative factor executes all of its selected source
calls, even if an earlier factor already makes the product zero.  At a layer $\ell\ge2$, its full arity $A_\ell$
has, uniformly over rows,
\[
 \mathbb E A_\ell\le1+3\rho^2,
 \qquad \mathbb E A_\ell^2\le1+18\rho^2,
 \qquad \rho\le r_{\ell-1}.
\]
For completeness, the pgf in the majority representation has
$\mathbb E K=\rho^2$ and
$\mathbb E K^2=\rho^2+\tfrac43\rho^4$ before the factor
$(\rho/\tanh\rho)/C_K$.  Since $C_K\ge1$ and
$\rho\coth\rho\le4/3$ for $\rho\le1$, these inequalities
follow by using $\sum_k\pi_k=1$ for the constant term:
\[
 \mathbb E_\pi(2K+1)\le1+\tfrac83\rho^2,\qquad
 \mathbb E_\pi(2K+1)^2
 =1+\mathbb E_\pi(4K^2+4K)
 \le1+\tfrac43(8\rho^2+\tfrac{16}3\rho^4)
 \le1+18\rho^2.
\]
The gate only reduces the arity.
Thus the mean and second moment are bounded by
$1+6/(\ell-1)$ and $1+36/(\ell-1)$.

Let $A_1$ be the resulting first-layer source arity, including the
zero arity when the outer gate closes.  The general-radius theorem
gives
\begin{equation}
 \mathbb E A_1\le2\max\{C,C^2\}\le2\kappa^2,\qquad
 \mathbb E[\text{local scalar work}]=O(\kappa^2B^4).
 \label{eq:column-growing-first-moment}
\end{equation}
For its second moment use the majority-index law
$\pi_k=(C/\tanh C)d_k/C_k$.  The proposal has
$\mathbb E_dK=C^2$ and
$\mathbb E_dK^2=C^2+\frac43C^4$.  Since
$\sum_k\pi_k=1$ and $C_k\ge1$,
\begin{align}
 \mathbb E A_1^2
 &\le \mathbb E_\pi(2K+1)^2 \nonumber\\
 &\le1+C\coth C\left(8C^2+\frac{16}3C^4\right)
 \le28\kappa^5.
 \label{eq:column-growing-second-moment}
\end{align}
The last step uses $C\coth C\le1+C\le2\kappa$.
The semantic first-layer count $S_1=1+A_1$ therefore satisfies
\[
 \mathbb E S_1\le3\kappa^2,\qquad
 \mathbb E S_1^2\le58\kappa^5.
\]
These bounds include the accepted full majority's source calls.
The scalar proposal and acceptance arithmetic is paid separately
by~\eqref{eq:column-growing-first-moment}.

For completeness, we give explicit bounds for the recurrence used
in the finite-bit argument.  Let $a_\ell,b_\ell$ bound the first two
moments of the total semantic vertices and source calls below a
normalized-neuron call in layer $\ell$.  The upper-layer arity
estimates give
\[
 a_\ell\le1+\mu_\ell a_{\ell-1},\qquad
 b_\ell\le\mu_\ell b_{\ell-1}
       +1+2\mu_\ell a_{\ell-1}+\nu_\ell a_{\ell-1}^2,
\]
where $\mu_\ell=1+6/(\ell-1)$ and
$\nu_\ell=1+36/(\ell-1)$.  The latter bounds the second factorial
arity moment as well.  Conditional on arity and predecessor
indices, fresh subcalls give this recurrence even when different
predecessor neurons have different costs.

Write $p_\ell=\prod_{j=1}^{\ell-1}(1+6/j)
=\binom{\ell+5}{6}\le(\ell+5)^6$.
Division of the first recurrence by $p_\ell$ yields
$a_\ell/p_\ell\le3\kappa^2+\ell-1$ and hence
\[
 a_\ell\le\kappa^2(\ell+5)^7.
\]
For $\ell\ge2$, the nonrecursive part of the second recurrence is
at most $52\kappa^4(\ell+4)^{14}$, since
$\mu_\ell\le7$ and $\nu_\ell\le37$.
Divide by $p_\ell$, bound each reciprocal by one, and sum.  This gives
\begin{equation}
 b_\ell\le p_\ell
  [58\kappa^5+52\kappa^4\ell(\ell+5)^{14}]
 \le110\kappa^5(\ell+5)^{21}.
 \label{eq:column-growing-tree}
\end{equation}

Each selected empirical coefficient uses at most $24D^2$
normalized or actual neuron roots, including the base function
coin when present.  If their complete recursive counts are $T_j$,
the pointwise inequality
$(\sum_{j=1}^{J}T_j)^2\le J\sum_{j=1}^{J}T_j^2$
does not require their independence.  With the linear diagram
bookkeeping included, a semantic count $N$ for a selected
coefficient has
\begin{equation}
 \mathbb E N^2\le
  2\cdot10^5\,\kappa^5(D+5)^{25}.
 \label{eq:column-growing-diagram}
\end{equation}
All arities, row choices, and monomial degrees use only known
weights and radii.  Full padding prevents source outcomes from
controlling these counts.  In particular the bound is uniform
over the empirical observations, the accepted split, and the
latent beta parameter.

We spell out the arithmetic charge to keep that last conditioning
statement separate from the implementation of the urn.  At the
$j$th reached primitive visit, the actual observable history
determines every rational parameter.  Fresh random bits implement
its next known coin.  The general-radius scalar theorem, the
unit-radius scalar bound at later layers, and exact rational
comparison give conditional expected local work
\[
 O\!\left(\kappa^2
  [B^2+\log(m+2)+\log(j+2)]^4\right).
\]
This statement concerns the observable filtration. It makes no
conditional precision-tail claim given the latent beta variable.
Summing the nonnegative predictable charges and using
$N\log^4(N+2)=O(N+N^2)$ with
\eqref{eq:column-growing-diagram} bounds the neuron and urn work.
The pairing, the finite Doob transform, and the four coordinate
sequences take $O(m[B^2+\log(m+2)]^4)$ further work.  Consequently
one selected coefficient costs at most
\begin{equation}
 O\!\left(
  \kappa^2[m+\kappa^5(D+5)^{25}]
       [B^2+\log(m+2)]^4\right)
 \label{eq:column-growing-coefficient}
\end{equation}
in expectation.  The constants in these bounds do not depend on
$C$ or on the point supplied to the sampler.

The correction levels have geometric ratio $1/4$ and
$m=2^am_0$.  Since $\log m_0=O(B)$, summing
\eqref{eq:column-growing-coefficient} over those levels gives
$O(\kappa^7(D+5)^{25}B^8)$ conditional work.  Here $\kappa^2$ is
the local scalar work of one visit, from
\eqref{eq:column-growing-first-moment}, and $\kappa^5$ enters
through $\mathbb E S_1^2$ in \eqref{eq:column-growing-diagram}.
The factor $B^8$ is the fourth power of the operand length $O(B^2)$.
The amplitude gate $q\le6D^{-1/2}$ and $32\le D\le B$ then give
$O(\kappa^7\sqrt D(D+5)^{24}B^8)=O(\kappa^7\sqrt D\,B^{32})$
online bits, including the initial branch selection.  This implies
the bound with $B^{40}$ stated in the theorem.

Finally, growing $C$ causes no preprocessing increase beyond the
original budget.  Each $c_j$ is a power of two between $2^{-b}$
and one; hence $C\le n$ has $O(b+\log n)$ bits.  The positive
order-four mass circuits remain of size $O(Dn^2)$ with operands
of $O(D[b+\log(nD+2)])=O(B^2)$ bits.  Schoolbook arithmetic uses
$O(Dn^2B^4)$ additional work and no larger order of storage.
At the stated precision this is $O(Dn^2\log^4n)$, within the
original $O(Dn^2\log^6n)$ budget.  The arbitrary-radius factory
creates its product proposal on demand and requires no table whose
size grows polynomially in $C$.  Substituting
$\kappa\le2K_{\mathrm{col}}$ gives the stated online bounds.
The mean-chain family has $C_{\mathrm{col}}=1$, so its critical
lower bound proves sharpness on each fixed bounded-envelope subclass.
This completes the proof of
Theorem~\ref{thm:column-envelope-critical}.

\paragraph{Scope and selection of the sampler.}
The column envelope is independent of algebraic rank.
Preprocessing may compare $K_{\mathrm{col}}^4$ with $D+1$ and
choose between this construction and the existing zero-bias sampler.
The resulting expected probe bound is
\[
 O\!\left(\min\{D+1,K_{\mathrm{col}}^2\sqrt{D+1}\}\right).
\]
Its bit work remains within the displayed conservative bound because
$D\le B$ in the stated regime.  The new source bound improves the
$O(D)$ envelope when $C_{\mathrm{col}}=o(D^{1/4})$.
Its bit bound is polylogarithmic whenever $C_{\mathrm{col}}$ is
polylogarithmic in $n$.  The exponent $40$ is an implementation
bound and makes no practical prediction.
